\chapter{Introduction}

Colorectal cancer (CRC) is one of the most significant malignancies worldwide, ranked as the third most common cancer diagnosis and a leading cause of cancer-related death \cite{bray2024globocan,iarcglobocan2022factsheet}. Improvements in surgery and adjuvant therapy have increased cure rates, but CRC patients remain at high risk of experiencing relapse. Approximately 20\%--30\% of patients with stage II--III colorectal cancer experience disease recurrence within 5 years after curative resection \cite{boute2024outcomes,nors2024recurrence}. Strikingly, this includes many patients who initially have favorable prognostic indicators for example, early-stage tumors with low TNM (Tumor--Node--Metastasis) stage (where T describes primary tumor extent, N regional lymph node involvement, and M distant metastasis) and strong lymphocytic infiltration. In clinical practice, such cases highlight the inadequacy of current stratification methods: classic methods of patient stratification may be insufficient to predict relapse accurately \cite{manzo2022}. As a result, there is an urgent need to identify additional biomarkers or features that can more accurately stratify CRC patients by recurrence risk.

Nowadays, scientific community is investigating the context in which the tumor resides, being aware that the tumor behaviour is dependent on the estabilishment of a certain tumor microenvironment (TME). CRC progression is influenced not only by the genetics of cancer cells but also by interactions with the immune system, stromal cells, and even the gut microbiota \cite{roelands2023atlas}. From these foundations, the MITICO (Microbiome Immune-System Tumor Interaction in Colorectal Oncology) research project was developed. It involves the recruitment of 240 early-stage CRC patients over five years, with the aim  to enable a more refined risk stratification that incorporates innovative biomarkers. This project involves the direct investigation of the composition and functional properties of the immune, stromal,  and microbiotic components within the tumors of the recruited patients. These biological data will then be integrated with clinical and lifestyle information for each patient in order to establish a robust and personalized method for patient stratification. \cite{mitico}. The underlying hypothesis is that novel risk factors such as specific immune cell phenotypes or microbial signals could complement standard clinicopathological metrics to better predict outcomes.

\section{Tumor-Associated Neutrophils (TANs) and CRC Prognosis}
\label{sec:TAN}
Among the various components of the CRC microenvironment, immune cells have attracted considerable interest as prognostic indicators and therapeutic targets. In particular, tumor-infiltrating lymphocytes (TILs) (such as cytotoxic CD8\textsuperscript{+} T cells) are well-known to correlate with improved survival in CRC, forming the basis of the ``Immunoscore'' approach to augment TNM staging \cite{pages2018immunoscore}. Nevertheless, other immune cell populations are not conventionally included in this type of assessment, despite the growing evidence that they can significantly influence both tumor behavior and immune response \cite{vanvlerkenysla2023myeloidstates,schmid2010myeloid}. Tumor-associated neutrophils (TANs) have emerged as a notable example of this complexity.

Neutrophils are the most abundant immune cells in the circulation and are classically the body's first responders to infection or injury \cite{wu2020}. They rapidly migrate from blood to tissue sites, where they can attack pathogens and orchestrate subsequent immune responses\cite{wu2020}. In cancer, their role is dual-edged and has been somewhat overlooked until recently \cite{wu2020}. Neutrophils do not show clearly a defined association with patient prognosis: on one hand they can exhibit anti-tumor functions by killing tumor cells, producing cytokines that recruit
and activate T cells, and constraining metastasis in certain contexts. \cite{wu2020}. On the other hand, neutrophils can acquire pro-tumoral functionality: they secrete growth factors that promote angiogenesis, facilitate tumor cell invasion, suppress T-cell activity (for instance through checkpoint ligand expression or cytokines like IL-8 and PGE$_{2}$), and form pre-metastatic niches \cite{nomenclature, complexrole}. This functional plasticity of neutrophils is often conceptualized as ``N1'' vs. ``N2'' polarization a spectrum between anti-tumor (N1) and pro-tumor (N2) phenotypes, analogous to M1/M2 macrophages \cite{wu2020}. In murine models and some human studies, interventions that skew neutrophils toward an N1 state (or deplete N2-like neutrophils) can inhibit tumor growth \cite{wu2020}. The TME can drive neutrophils polarization toward N2 functional state, usually in a context in which TGF-$\beta$  is the predominant signal. N2 neutrophils are related with tumor immune evasion and cancer progression \cite{wu2020}.

\textbf{Prognostic ambiguity:} Because neutrophils can play these opposing roles, the mere presence or quantity of TANs in a tumor has not yielded a consistent prognostic pattern in CRC. Some studies reported that high neutrophil counts correlate with poor outcomes (supporting the notion of pro-tumor N2 dominance), while others found no correlation or even an association with better outcomes in certain settings \cite{manzo2022}. A consensus is emerging that neutrophil phenotype (i.e. the functional state of secreted molecules and the downstream effects on other cells), rather than just abundance, is the determinant factor that shape their role in tumor progression or tumor control \cite{que2022neutrophiltherapies,wu2024tans}. Recent comprehensive reviews stress neutrophil diversity and plasticity in cancer, noting that neutrophils can shift functions depending on signals in the TME \cite{manzo2022}. Therefore, identifying specific neutrophil subpopulations and biomarkers is key for leveraging TANs as prognostic indicators \cite{manzo2022}.

\subsection{A Distinct Neutrophil Subset (CD15\textsuperscript{high} TANs)}
\label{sec:IEO}

A breakthrough in this regard came from a 2022 study by Manzo et al., which performed high-dimensional immune profiling on CRC tumors \cite{manzo2022}. The researchers discovered that CRC tumors could be stratified based on the presence of a particular neutrophil subset defined by high expression of CD15. CD15, also known as Lewis X antigen, is a carbohydrate adhesion molecule present on neutrophils from early stages of their development \cite{nomenclature}. Functionally, CD15 acts as a ligand for E-selectin, mediating neutrophil tethering and rolling on endothelial cells during inflammation. Perhaps more importantly, engagement of CD15 can results in a signal cascade in neutrophils: experiments have shown that upon functional interaction of CD15 on leukocyte, it surfaces triggers cell activation and release of inflammatory cytokines \cite{lo1997}. Thus, CD15 is not only a marker of neutrophil differentiation but also a participant in neutrophil priming and functional activation.

In the study by Manzo et al., investigators noted that some CRC patients shows an engraftment of neutrophils expressing high levels of CD15 within the tumor, while other patients exhibit a reduced amount of CD15high neutrophils infiltrating the tumor. They defined two groups of patients accordingly: HN (High CD15\textsuperscript{high} Neutrophils) and LN (Low CD15\textsuperscript{high} Neutrophils), using a 50\% threshold (i.e. HN patients have $>$50\% of intratumoral neutrophils that are CD15\textsuperscript{high}, whereas LN patients have $<$50\%) \cite{manzo2022}. Importantly, this stratification revealed strong associations with immune context and outcomes:

\begin{itemize}
    \item \textbf{Immune context:} HN tumors were found to harbor an ``atypical'' immune activation profile. Specifically, they were enriched in CD8\textsuperscript{+} T cells that produced Granzyme K (GZMK) instead of the conventional Granzyme B. This GZMK\textsuperscript{high} CD8\textsuperscript{+} effector memory T-cell population was rarely seen in LN tumors \cite{manzo2022}. The presence of GZMK\textsuperscript{high} T cells indicates a skewed T-cell response, which is not tumor specific and can led to chronic injury of the tissue. Mechanistic studies suggested that neutrophils in HN tumors, through contact-dependent signals and cytokines like CXCL12, were driving this T-cell skewing \cite{manzo2022}. In essence, CD15\textsuperscript{high} neutrophils were interacting with T cells and dampening their effective anti-tumor function, even though the number of T cells stay high in the tumor and this can be misleading in the risk assessment of the patient.
    
    \item \textbf{Clinical outcome:} Patients in the HN group had significantly worse clinical outcomes than those in the LN group \cite{manzo2022}. In non-metastatic CRC, HN status correlated with higher rates of recurrence and poorer disease-free survival. This suggests that CD15\textsuperscript{high} neutrophils are a marker of aggressive tumor behavior. The study's multivariate analysis indicated that the HN/LN classification could add prognostic value on top of traditional factors. In other words, even among patients with similar stage and lymphocyte infiltration, those with an abundance of CD15\textsuperscript{high} neutrophils have a poorer prognosis, highlighting the pro-tumor influence of this neutrophil subset \cite{manzo2022}.
\end{itemize}

Collectively, these findings introduced CD15\textsuperscript{high} TANs as a candidate biomarker for CRC stratification. The biological model is that CD15\textsuperscript{high} neutrophils, retained in the tumor by chemokines (e.g. CXCL12) \cite{manzo2022}, establish an immunosuppressive milieu by skewing T cells and possibly other immune cells, thereby promoting tumor progression \cite{manzo2022}. This provided a compelling rationale for the clinicians and engineers in our team to try to detect this phenomenon in a more scalable way.

\subsection{Limitations of Current TAN Profiling Methods}
\label{sec:limitations}
While the HN vs. LN classification is promising, determining a patient's status currently requires specialized laboratory procedures. Tumor tissue must be freshly collected and processed to yield single-cell suspensions for flow cytometry (or subjected to multiplex immunofluorescence). The analysis involves staining for multiple markers (e.g. CD66b to identify neutrophils, CD15, CD10, HLA-DR, etc.) \cite{manzo2022} and often also functional assays or imaging. These steps are resource-demanding since they need important amount of tissue material that might not be available outside research settings, expensive reagents and instruments, and skilled personnel to perform and interpret the assays. Moreover, the workflow is inherently slower and separate from the standard diagnostic process, which relies mainly on H\&E-stained slides. From a practical standpoint, this means HN/LN stratification is not yet ready for routine clinical use despite its potential value. It may also miss macroscopic histopathological patterns: for instance, certain architectural features or stromal configurations in H\&E slides could correlate with immune infiltration patterns, but such correlations would be overlooked by focusing only on the phenotype of single cells obtained from the tumor processing.

\textbf{Opportunity for AI:} These limitations lead to the outstanding question : can we circumvent the laborious workflow by using the available pathology slides? If an AI algorithm can learn to recognize the morphological footprints of an ``HN-type'' tumor (for example, abundant neutrophils in tissue sections, or perhaps more subtle signs like tissue damage or specific stromal features), it could provide a fast, cost-effective proxy for the CD15\textsuperscript{high} neutrophil biomarker. This would effectively democratize the stratification: any hospital with a digital slide scanner could potentially apply the model, without need for new wet-lab tests. Additionally, this approach might capture a more comprehensive picture since the model could, in principle, integrate neutrophil presence with other histological cues (like the spatial distribution of lymphocytes, necrotic areas, mucin presence, lipid composition etc.) that might also relate to prognosis.

\section{AI in Digital Pathology for Immunology}
\label{sec:AIinDigPat}
Digital pathology refers to the use of high-resolution whole-slide images (WSIs) for routine diagnostics and quantitative analysis, enabling computational methods to support pathologists and large-scale biomarker discovery \cite{tizhoosh}. Over the past decade, deep learning has increasingly been applied to WSIs for tasks such as tumor detection, grading, and prediction of clinically relevant biomarkers from routine H\&E slides \cite{tizhoosh} \cite{asif} . In terms of clinical translation, a milestone was reached with the FDA De Novo authorization of an assistive AI system for digital pathology (Paige Prostate; DEN200080), highlighting the growing regulatory and clinical interest in these tools \cite{fda}. Beyond headline performance, generalization and reliability are essential when proposing AI models for clinical use. We therefore enforce strict patient-level separation: tiles/patches from the same patient never appear in both training and validation sets, preventing subtle data leakage that can occur in tiling-based pipelines \cite{bussola}. Such leakage can substantially inflate performance when near-duplicate tissue regions from the same slide or patient are inadvertently split across training/validation/test \cite{bussola}. 
% In addition, we explicitly assess domain shift by evaluating the final models on an independent external cohort, since differences in staining, scanning, and patient populations can degrade performance and distort probability calibration outside the training distribution \cite{asif}. 
Finally, we consider confidence and calibration as first-class aspects of evaluation, since deep networks can be overconfident—especially under distribution shift—and uncertainty-aware strategies can support safer high-confidence predictions \cite{dolezal}.
For research applications, numerous studies have shown artificial intelligence (AI) models' ability to extract prognostic and predictive biomarkers from routine H\&E slides:

\begin{itemize}
    \item \textbf{Molecular surrogates:} Deep models have inferred genotype and phenotype information from H\&E. Examples include prediction of specific gene mutations, or transcriptomic signatures directly from tumor histology. These works suggest that subtle patterns in cell morphology and tissue architecture often imperceptible to human observers - can correlate with underlying molecular alterations, and AI can exploit this.
    
    \item \textbf{Immunophenotypes:} Some studies focus on immune-related labels.A representative example used weakly supervised learning to predict p16 status (a surrogate for HPV infection in oropharyngeal cancer) directly from H\&E WSIs, and provided interpretable evidence by highlighting histological regions associated with p16 positivity \cite{adachi}. Similarly, efforts to predict the density of tumor-infiltrating lymphocytes or the presence of tertiary lymphoid structures from H\&E have shown that AI can quantify immune infiltration with reasonable concordance to manual scoring \cite{manzo2022, nomenclature}.
    
    \item \textbf{Outcomes and risk scores:} Beyond specific biomarkers, there are AI models that directly predict patient outcomes (survival, recurrence) from H\&E images \cite{pawlik2025}. These models effectively learn their own prognostic features, which sometimes overlap with known factors (like mitotic count or lymphocytic infiltration) but can also include novel image-based patterns.
\end{itemize}

Given this context, applying Al to predict the neutrophil-driven stratification (HN vs. LN) is a natural extension. It represents a specific case of in silico immune profiling, where the aim is to have the computer replicate what a specialized lab test would tell us about the tumor's immune status.

\subsection{Weakly Supervised Learning with Multiple Instance Learning (MIL)}
\label{sec:MIL}
Whole-slide images pose unique challenges due to their size (often over 1 gigapixel per slide) and heterogeneity. A single WSI can contain various tissue components tumor, stroma, normal mucosa, immune infiltrates, necrosis only some of which may be relevant for predicting a given label. Moreover, in our case the label (HN or LN) is a property of the entire slide/tumor, not of any specific location, and we lack localized annotations of neutrophil-rich regions, and thus rely on weak supervision with slide-level labels \cite{campanella2019}. To tackle this, our approach uses a Multiple Instance Learning (MIL) framework \cite{Lu,campanella2019}.

In MIL, each slide is divided into many smaller patches (also called tiles or instances). These patches (e.g. $256\times256$ pixel regions at high magnification) are treated as individual inputs to a neural network that extracts features from them \cite{Lu,campanella2019}. The slide-level prediction is then an aggregation of all patch-level features. Conceptually, one can imagine the model scanning the slide patches to find evidence of either class: for an HN slide, perhaps certain patches contain dense clusters of neutrophils or specific tissue damage patterns that the model can learn to recognize; for an LN slide, such patches might be absent or different in character. The model's job is to identify and weight the informative patches to make an overall prediction.

\textbf{Attention-based MIL:} A critical component we employ is an attention mechanism within the MIL model \cite{ilse2018}. Rather than assuming that all patches contribute equally, or that only the ``maximum prediction'' patch matters, attention allows the model to learn how much each patch should contribute. The network produces both a feature embedding for each patch and an attention weight. During training, it optimizes these weights such that patches predictive of the slide's label are given higher weight. This not only improves performance but also yields interpretability: by inspecting the learned attention weights on a new slide, we can visualize which regions the model deemed important for its decision. We use the attention-MIL approach originally described by \cite{ilse2018}. and further refined in methods like Clustering-constrained Attention Multiple Instance Learning (CLAM), which incorporates clustering to help discover subtypes of patches \cite{Lu}. These approaches are well-suited to our problem, as we expect only a minority of regions in an HN tumor, for instance, to show the neutrophil-related pattern, and attention can focus the model on those.

\textbf{Transformer MIL:} We also experiment with a Transformer-based MIL model (inspired by TransMIL \cite{shao2021}) that uses self-attention across all patches. This allows the model to consider patch-patch interactions (e.g. a patch of tumor epithelium might be interpreted differently if an adjacent patch has a lymphoid aggregate). While more computationally intensive, the transformer MIL may capture contextual cues indicative of HN status, such as spatial proximity of neutrophils to tumor cells or overall inflammatory patterns.

\subsection{Leveraging Pretrained ``Foundation'' Models}
\label{sec:foundation}
A major challenge in training MIL models on pathology data is the required dataset size. Training a deep CNN from scratch on tens of thousands of patches (per slide) across hundreds of slides can easily involve billions of pixels and risk overfitting if the dataset is small (in our case, we have on the order of 100 slides). To address this, we harness transfer learning via self-supervised pretrained models. Specifically, we use a Vision Transformer (ViT) model that was pretrained on a large corpus of histology images using a self-supervised method (Masked Image Modeling). This model - akin to recent pathology foundation models like CLAMOR or Virchow - has learned rich representations of histological structures by looking at myriad examples from varied tissue types \cite{virchow}. We treat this ViT as a fixed feature extractor: each WSI patch is fed through the ViT to obtain an embedding (a vector in a latent feature space). These embeddings capture high-level information about the patch such as whether it contains lymphocytes, necrosis, glandular structures, etc., as the model has presumably learned those during pretraining. We then train our MIL classifier on these embeddings instead of raw images.

This brings two key advantages:
\begin{itemize}
    \item \textbf{Data efficiency:} The pretrained features are more informative and generalizable, allowing our MIL model to learn the HN vs. LN discrimination with far fewer slides than would otherwise be needed. Prior work has shown that using such embeddings can significantly boost accuracy on small pathology datasets. In our experiments, this was crucial to achieve good performance.
    \item \textbf{Computational efficiency:} We avoid back-propagating gradients through a massive CNN for every patch in every epoch. Instead, we only train the relatively smaller MIL aggregator network, while the heavy image processing is done once by the fixed backbone. We further accelerate training by caching all patch embeddings for each slide, so that the model can iterate quickly over slides during learning. This setup is amenable to HPC parallelization : we computed embeddings for all slides in parallel across many nodes.
\end{itemize}

\section{Problem Statement}
\label{sec:problem}

\textbf{Core question:} Can an Al model predict whether a CRC tumor is ``HN'' or ``LN'' (high vs. low CD15\textsuperscript{high} neutrophils) using only the H\&E whole-slide image? This problem statement encapsulates a specific instance of computational pathology for immunology. Answering it involves several sub-questions:
\begin{itemize}
    \item Do H\&E images contain visual correlates of neutrophil infiltration and phenotypes? For example, heavy neutrophil presence might manifest as certain textural patterns or subtle color differences (since neutrophils have characteristic morphology on H\&E: multi-lobed nuclei and pale cytoplasm). If present, can a model learn to detect these reliably?
    \item Even if neutrophils are not individually detectable (H\&E cannot differentiate neutrophils from other granulocytes easily, especially without special stains), perhaps there are indirect clues. For instance, neutrophil-rich tumors might have more tissue necrosis (as neutrophils can cause collateral damage) or certain stromal changes (edema, fibrosis patterns) that a model could pick up on. Essentially, the model might learn a proxy representation of the immune microenvironment.
    \item What level of accuracy can be achieved, and is it sufficient to be useful? We aim to compare the model's predictions to the ground truth from flow cytometry. A successful model would have high slide-level classification performance (AUC, accuracy) for HN vs. LN.
    % \item How generalizable is the model? We must ensure it's not overfitting to cohort-specific artifacts (batch effects, scanner differences, etc.). This is why we plan external validation.
    \item Can the model provide explanations for its predictions that align with biological expectations (e.g. highlighting regions with heavy inflammation)? This is important for trust in a clinical context.
    \item If successful, this approach could serve as a prototype for predicting other immune markers from H\&E (for example, trying to predict ``Immunoscore'' or presence of specific immune cell types from morphology). Thus, a broader objective is to advance methodologies for virtual immune profiling.
\end{itemize}

\textbf{Aim:} The aim of this thesis is focused on developing and evaluating a deep learning model for the specific prediction of neutrophil-associated risk in CRC. We utilized a retrospective CRC cohort that has ground truth HN/LN labels from the MITICO immune analyzes. The dataset is relatively small (dozens of HN and LN cases), which led our choice to use pretrained features. We emphasize methodological rigor in model training (cross-validation, avoiding leakage) and include a test on 13 held-out patients for generalization. While our analysis touches on clinical outcome correlation (exploring if model-predicted HN vs. LN relates to patient recurrence outcomes), we treat that as a secondary, hypothesis-generating analysis. The primary goal is the feasibility demonstration of predicting an immunological label from histology. Detailed optimization of the model for deployment or its integration into clinical decision pipelines is beyond our current investigation.
We also note that our work is limited to a single-center training dataset; thus, conclusions about generality must be tempered by the need for broader validation.

\textbf{Thesis Contributions and Outline:} This project combines knowledge from oncology, immunology, and machine learning. The main contributions are summarized below, and each corresponds to sections of the thesis:
\begin{itemize}
    \item \textbf{Contribution 1:} A novel application of attention-based and transformer-based MIL for predicting a complex immunological phenotype (neutrophil polarization) from WSIs. To our knowledge, this is the first work to target tumor-associated neutrophil phenotypes via H\&E image analysis in CRC. This contributes to the emerging field of computational pathology for immune profiling.
    \item \textbf{Contribution 2:} Integration of foundation model features and HPC-optimized pipeline design. By using transformer-derived histology embeddings and scalable computing techniques, we demonstrate an efficient workflow that can handle gigapixel images reproducibly. This engineering contribution is important for moving such AI models toward practical utility, ensuring that results can be generated and verified on large datasets.
    \item \textbf{Contribution 3:} Interpretability-driven validation and clinical relevance analysis on a held-out internal cohort. 
    Beyond reporting aggregate performance, we provide slide-level interpretability by generating attention-based heatmaps and overlays that highlight the tissue regions driving the model’s predictions; these visualizations are intended for expert review performed by pathologists, in order to assess whether the model focuses on plausible histomorphological patterns. 
    We additionally evaluate probability calibration to ensure that predicted HN/LN risks can be interpreted meaningfully, rather than as uncalibrated scores. 
    Finally, we explore the clinical relevance of the model outputs by analyzing how the predicted risk probabilities correlate with patient Disease Free Survival, bridging the computational model back to the prognostic application of identifying patients at high risk of relapse. 
    Given the limited sample size, these analyses are presented as exploratory and hypothesis-generating rather than definitive.

\end{itemize}

The thesis is organized as follows. Chapter 2 describes the materials and dataset, including patient cohorts and labels, slide digitization and preprocessing, and the tile extraction and quality-control pipeline used to build the learning dataset. Chapter 3 presents the proposed methods, detailing the feature extraction strategy based on pretrained foundation models and the Multiple Instance Learning (MIL) architectures, together with the training protocol and the evaluation metrics. Chapter 4 reports the experimental results, including cross-validation performance and evaluation on a held-out internal test cohort, as well as interpretability analyses based on attention heatmaps and representative case studies. In the same chapter, we investigate the clinical relevance of the model outputs on an external cohort (TCGA) lacking CD15-derived ground-truth labels, by performing label-free inference and assessing whether the predicted HN-risk probabilities are associated with patient Disease Free Survival.
Chapter 5 discusses the implications of the findings in the context of related work, highlights limitations and outlines future directions such as multimodal integration and prospective validation. Finally, Chapter 6 concludes the thesis by summarizing the main contributions and their potential impact on computational pathology for colorectal cancer.

